\psection[1.0]{background}{0.85}{I1,I2,I4}{core,dal}{Background and Motivation}
% Section II. Literature only. Every claim was checked on 16 Sep 2026 against
% the verified bib notes and the cited papers' abstracts or text (COMPL-AI's
% simple-average rule: arXiv 2410.07959, aggregation paragraph; Biswas and
% Rajan: Finding 5; Kozodoi et al.: abstract; Wachter et al.: abstract).
% Scenarios follow the writing brief: reduced merges the toolkit and
% SE-fairness paragraphs; withdrawn cuts them to two sentences and makes the
% validity of the language-model fairness instruments central.

\textbf{Regulation and its technical interpretation.}
The AI Act attaches obligations to high-risk systems, among them data governance (Art.~10) and accuracy, robustness and cybersecurity (Art.~15), without saying how any of them is to be measured~\cite{euaiact2024}.
Its principles draw on the non-binding guidelines of the High-Level Expert Group on AI~\cite{hleg2019trustworthyai}.
The NIST AI Risk Management Framework and ISO/IEC~42001 organise risk management across an organisation and prescribe no test for an individual model~\cite{nist2023airmf,iso42001_2023}.
Requirements-engineering research extracts and tracks legal requirements, the AI Act included, and automates compliance checking~\cite{abualhaija2025llmregextraction,abualhaija2024regchange,amaral2022nlpcompliance}.
COMPL-AI is the closest technical interpretation of the Act for language models~\cite{guldimann2024complai}.
Its eighteen technical requirements fall under six principle groups, and twelve of them are measurable with public benchmarks.
It combines scores by a simple average at every level, on the stated ground that the regulation sets no hierarchy among requirements.
VERA takes over the twelve measurable requirements, \req{01} to \req{12}, as its unit of review.

\textbf{Measurement and its validity.}
Evaluation suites for language models reach beyond accuracy.
HELM reports calibration, fairness and toxicity over shared scenarios~\cite{liang2023helm}, and DecodingTrust covers perspectives such as stereotype bias and privacy~\cite{wang2023decodingtrust}.
The language model evaluation harness runs a large task library behind one interface~\cite{gao2021lmevalharness}, garak scans models for security vulnerabilities~\cite{derczynski2024garak}, and PromptOps applies metamorphic testing to robustness and fairness~\cite{sunetnanta2025promptops}.
These tools settle what is measured.
How the measurements combine into a judgement is fixed inside the suite or left to the reader.
In neither case is the combination rule an input that the reviewing institution declares, versions and answers for.

A score is also only as good as its construct.
Raji et al.\ argue that influential benchmarks are presented as measures of general progress that their construction cannot support~\cite{raji2021everything}, and BetterBench finds that most of the benchmarks it assessed report no statistical significance and are hard to replicate~\cite{reuel2024betterbench}.
Bias benchmarks for language models are a documented weak point: an inventory of stereotyping benchmarks, StereoSet~\cite{nadeem2021stereoset} among them, found that they often do not articulate what they measure~\cite{blodgett2021stereotyping}.
VERA's language-model specification scores bias on StereoSet, BBQ~\cite{parrish2022bbq} and BOLD~\cite{dhamala2021bold}, and the last two were outside that inventory.
We therefore do not use language-model bias scores as fairness evidence: \Cref{sec:fairness} applies fairness criteria to decisions with outcome labels instead.

\ifarmfull
\textbf{Fairness criteria and their choice.}
For tabular classifiers, AIF360, Fairlearn and Aequitas make fairness criteria easy to compute~\cite{bellamy2019aif360,weerts2023fairlearn,saleiro2018aequitas}.
The criteria formalise different notions of fairness, grouped as independence, separation and sufficiency, and the same classifier can be fair under one and unfair under another~\cite{verma2018fairness,barocas2023fairness}.
The toolkits leave the choice among criteria, and its justification, to the user.
Practitioners report needs that the toolkits do not meet, such as help in placing fairness in context and communicating it to stakeholders~\cite{lee2021landscape,deng2022exploring}.

Software engineering treats fairness as a testable quality attribute~\cite{brun2018software}, and Themis generates discrimination tests without an oracle~\cite{galhotra2017fairness}.
Adult~\cite{becker1996adult} and German Credit~\cite{hofmann1994statlog} are among the public datasets of fairness-testing research~\cite{chen2024fairness,chakraborty2021bias}, and the tabular specification of \Cref{sec:vera} draws on both\decision{D6}{legal,core}{22 Sep}{German Credit version (UCI Statlog coding or South German Credit, gromping2019south) and protected attribute; the dataset citation in Section II follows the answer}.
On Kaggle models that include credit-risk tasks, Biswas and Rajan find that a model can look fair on disparate impact and worse on error rates~\cite{biswas2020machine}.
Bias mitigation is a separate field~\cite{hort2024bias}; VERA measures and does not mitigate.
\fi

\ifarmreduced
\textbf{Fairness criteria and their choice.}
Toolkits such as AIF360, Fairlearn and Aequitas compute fairness criteria for tabular classifiers~\cite{bellamy2019aif360,weerts2023fairlearn,saleiro2018aequitas}.
The criteria formalise different notions of fairness, so the same classifier can be fair under one and unfair under another~\cite{verma2018fairness,barocas2023fairness}.
The toolkits leave the choice, and its justification, to the user, while practitioners ask for help placing fairness in context~\cite{lee2021landscape,deng2022exploring}.
Software engineering tests fairness as a quality attribute~\cite{brun2018software,galhotra2017fairness,chakraborty2021bias}, and different metrics rank the same models differently~\cite{biswas2020machine}.
A survey of the field lists Adult~\cite{becker1996adult} and German Credit~\cite{hofmann1994statlog} among its datasets~\cite{chen2024fairness}, and the tabular specification of \Cref{sec:vera} uses both\decision{D6}{legal,core}{22 Sep}{German Credit version (UCI Statlog coding or South German Credit, gromping2019south) and protected attribute; the dataset citation in Section II follows the answer}.
VERA measures bias and does not mitigate it~\cite{hort2024bias}.
\fi

\ifarmwithdrawn
\textbf{Fairness criteria and their choice.}
Fairness toolkits compute many criteria~\cite{bellamy2019aif360,weerts2023fairlearn}, a classifier can be fair under one and unfair under another~\cite{verma2018fairness,barocas2023fairness}, and practitioners ask for help placing fairness in context~\cite{lee2021landscape,deng2022exploring}.
Software engineering tests fairness as a quality attribute~\cite{brun2018software,galhotra2017fairness,chen2024fairness}, and different metrics rank the same models differently~\cite{biswas2020machine}.
\fi

In credit scoring, the choice of criterion has been studied directly.
Kozodoi et al.\ recommend separation for scorecards and find on real data that several criteria can be approximately satisfied at once~\cite{kozodoi2022fairness}.
Hurlin et al.\ show how lenders and their regulators can test the fairness of a scoring model formally and identify the variables behind a lack of fairness~\cite{hurlin2026fairness}.
Wachter et al.\ argue that EU non-discrimination law is too contextual to be automated, and they propose conditional demographic disparity as a baseline statistic whose interpretation stays with the courts~\cite{wachter2021fairness}.
Practitioner studies document gaps between fairness research and the challenges of product teams~\cite{holstein2019improving}, and fairness work that rests on ad hoc processes and individual advocates~\cite{madaio2020codesigning}.
In a bank, the choice of criterion is also open to challenge by a control function that did not build the model (\Cref{sec:context}).
No tool discussed above records the scoring policy, fairness criterion included, as a versioned review artefact that one reviewer can cite, another can challenge, and a later run can be checked against.
This paper treats that policy as data (I1), makes the fairness criterion an explicit entry in it (I2), and reports what introducing it into a bank's review process involved (I4).
